% Options for packages loaded elsewhere
\PassOptionsToPackage{unicode}{hyperref}
\PassOptionsToPackage{hyphens}{url}
%
\documentclass[
  11pt,
]{article}
\usepackage{amsmath,amssymb}
\usepackage{lmodern}
\usepackage{setspace}
\usepackage{iftex}
\ifPDFTeX
  \usepackage[T1]{fontenc}
  \usepackage[utf8]{inputenc}
  \usepackage{textcomp} % provide euro and other symbols
\else % if luatex or xetex
  \usepackage{unicode-math}
  \defaultfontfeatures{Scale=MatchLowercase}
  \defaultfontfeatures[\rmfamily]{Ligatures=TeX,Scale=1}
\fi
% Use upquote if available, for straight quotes in verbatim environments
\IfFileExists{upquote.sty}{\usepackage{upquote}}{}
\IfFileExists{microtype.sty}{% use microtype if available
  \usepackage[]{microtype}
  \UseMicrotypeSet[protrusion]{basicmath} % disable protrusion for tt fonts
}{}
\makeatletter
\@ifundefined{KOMAClassName}{% if non-KOMA class
  \IfFileExists{parskip.sty}{%
    \usepackage{parskip}
  }{% else
    \setlength{\parindent}{0pt}
    \setlength{\parskip}{6pt plus 2pt minus 1pt}}
}{% if KOMA class
  \KOMAoptions{parskip=half}}
\makeatother
\usepackage{xcolor}
\usepackage[margin = 1in]{geometry}
\usepackage{graphicx}
\makeatletter
\def\maxwidth{\ifdim\Gin@nat@width>\linewidth\linewidth\else\Gin@nat@width\fi}
\def\maxheight{\ifdim\Gin@nat@height>\textheight\textheight\else\Gin@nat@height\fi}
\makeatother
% Scale images if necessary, so that they will not overflow the page
% margins by default, and it is still possible to overwrite the defaults
% using explicit options in \includegraphics[width, height, ...]{}
\setkeys{Gin}{width=\maxwidth,height=\maxheight,keepaspectratio}
% Set default figure placement to htbp
\makeatletter
\def\fps@figure{htbp}
\makeatother
\setlength{\emergencystretch}{3em} % prevent overfull lines
\providecommand{\tightlist}{%
  \setlength{\itemsep}{0pt}\setlength{\parskip}{0pt}}
\setcounter{secnumdepth}{-\maxdimen} % remove section numbering
\usepackage{amsmath}
\usepackage{bbm}
\usepackage{array}
\usepackage{multirow}
\usepackage{graphicx}
\usepackage{float}
\usepackage{apacite}
\usepackage{natbib}
\ifLuaTeX
  \usepackage{selnolig}  % disable illegal ligatures
\fi
\IfFileExists{bookmark.sty}{\usepackage{bookmark}}{\usepackage{hyperref}}
\IfFileExists{xurl.sty}{\usepackage{xurl}}{} % add URL line breaks if available
\urlstyle{same} % disable monospaced font for URLs
\hypersetup{
  pdfauthor={Seyed Mohammad Mehdi Hassani Najafabadi(Student ID: 400489126)},
  hidelinks,
  pdfcreator={LaTeX via pandoc}}

\title{\includegraphics[width=1in,height=\textheight]{Logo.png}

STATS/CSE 780

Assignment3 1}
\author{Seyed Mohammad Mehdi Hassani Najafabadi(Student ID: 400489126)}
\date{27 March, 2023}

\begin{document}
\maketitle

{
\setcounter{tocdepth}{2}
\tableofcontents
}
\setstretch{1.5}
\newpage

\hypertarget{introduction}{%
\section{Introduction}\label{introduction}}

The following
dataset\href{https://archive.ics.uci.edu/ml/datasets/bank+marketing}{{[}1{]}}
aims to develop a customer segmentation analysis to define an effective
marketing strategy for a credit card company. The data set used in this
analysis summarizes the usage behavior of approximately 9000 active
credit card holders during the past six months. The dataset is available
at the customer level and comprises 18 behavioral variables that provide
valuable insights into customer behavior. The behavioral variables in
the dataset include factors such as transaction frequency, transaction
amount, balance, credit limit, and payment amount. Other variables
include demographic information such as age, gender, and income, as well
as variables that describe customer behavior, such as tenure, loyalty,
and credit card type. The customer segmentation analysis will help the
credit card company divide its customer base into distinct groups with
similar characteristics, needs, and behaviors. This analysis will
provide insights into the different groups' preferences, enabling the
company to develop a targeted marketing strategy that is tailored to
each group's specific needs.

\hypertarget{data-transformation-and-preprocessing}{%
\section{Data Transformation and
Preprocessing}\label{data-transformation-and-preprocessing}}

\hypertarget{handling-missing-values}{%
\subsection{Handling missing values}\label{handling-missing-values}}

There are 1690 missing values(NA) which will be removed from the data
set:

\hypertarget{outliers-analysis}{%
\subsection{Outliers analysis}\label{outliers-analysis}}

According to the box plot in the supplementary document, there are some
outliers which can be handled by IQR(Inter Quartaile Range) technique:

In order to select important variables,pricipal components analysis is
used as follows:

\hypertarget{cluster-analysis}{%
\section{Cluster Analysis}\label{cluster-analysis}}

\hypertarget{k-means-clustering}{%
\subsection{K-means Clustering}\label{k-means-clustering}}

\hypertarget{determine-the-number-of-clusters}{%
\subsubsection{Determine the number of
clusters}\label{determine-the-number-of-clusters}}

In order to use effective clustering, firstly the optimal number of
clusters has been extracted by using Elbow method width method as
follows:

\hypertarget{elbow-method}{%
\subsubsection{Elbow method}\label{elbow-method}}

According to the Elbow and Silhouette the optimal number of clusters is
K=3.In the following section, k=3 considered as a cluster numbers.

\hypertarget{get-average-silhouette-width}{%
\subsubsection{Get Average silhouette
width}\label{get-average-silhouette-width}}

Seems relative good clustering with kmeans.

\hypertarget{hierarchical-clustering}{%
\subsection{Hierarchical Clustering}\label{hierarchical-clustering}}

In order to use effective Hierarchical Clustering clustering, firstly
the optimal number of clusters has been extracted by using Elbow method
width method as follows

\hypertarget{elbow-method-for-hiercahical-clustering}{%
\subsubsection{Elbow method for hiercahical
clustering}\label{elbow-method-for-hiercahical-clustering}}

the best cutoff in this clustering is 4.

Again the results of hierarchical is fairly good. clustering with

\hypertarget{pricipal-components-analysis}{%
\subsection{Pricipal Components
Analysis}\label{pricipal-components-analysis}}

Firstly the number of principle components are extracted by Elbow
Methods: As It can be seen two principle components is acceptableness
for the next combination by other techniques:

\hypertarget{pricipal-components-analysis-and-hierarchical}{%
\subsection{Pricipal Components analysis and
Hierarchical}\label{pricipal-components-analysis-and-hierarchical}}

By Combining Principal Components analysis and Hierarchical, silhouette
score is score remain stable.

\hypertarget{visualize-clustering-results-of-using-pca-and-kmeans}{%
\section{Visualize Clustering Results of using PCA and
KMeans}\label{visualize-clustering-results-of-using-pca-and-kmeans}}

After applying PCA, the quality of K-Means has been dramatically
increased to 0.53 Silhouette average score

\hypertarget{comparison-with-different-clustering-algorithms}{%
\section{Comparison with different clustering
algorithms}\label{comparison-with-different-clustering-algorithms}}

Comparison 1 : According to the obtained silhouette average between
K-means and hierarchical, hierarchical clustering has better results
(Silhouette average:0.56) than K-Means with (Silhouette average:0.56).

Comparison 2: After applying PCA, the quality of K-Means has been
dramatically increased to 0.53 Silhouette average score. However,the
mixture of Hierarchical and PCA performance remains stable.

\newpage

\hypertarget{references}{%
\section{References}\label{references}}

\href{https://archive.ics.uci.edu/ml/datasets/bank+marketing}{1}.
\url{https://archive.ics.uci.edu/ml/datasets/bank+marketing}

\end{document}
